%
% menu_lifecycle_d2_buggy.tex
%
% FIDELITY CONTROL for docs/menu_lifecycle.tex, defect D2.
%
% This spec differs from the fixed menu_lifecycle.tex in exactly one place:
% DepartLapsed does NOT sweep menuOwner. This is the shipped wiring -- the
% DepartureCascade fires the inbox sink on the lapse path but the menu
% registry prune is wired only to session_cleanup.py, which the background
% timer never runs. A timer-reaped owner leaves the registry with its bar
% entry still in menuOwner: an orphaned-owner menu, MO's violation.
%
% Expected: MO's negation is reachable.
%   probcli docs/menu_lifecycle_d2_buggy.tex -model_check -nodead \
%     -goal "#c.(c : CONN & c : menuOwner & not(c : registered))" \
%     -p DEFAULT_SETSIZE 3
%   => goal FOUND (Connect; Identify; SetMenu; LeaseLapse; DepartLapsed).
%
% Against the fixed spec the same goal is NOT found. This control isolates
% D2: SetMenu is atomic (D3 fixed) and DepartGraceful still sweeps the menu,
% so the only reachable violation is the lapsed departure path.
%

\documentclass[a4paper,10pt,fleqn]{article}
\usepackage[margin=1in]{geometry}
\usepackage{fuzz}
\usepackage[colorlinks=true,linkcolor=blue,citecolor=blue,urlcolor=blue]{hyperref}

\begin{document}

\title{lux Agent-Menu Lifecycle --- D2 Fidelity Control}
\author{The lapse path that reaps the inbox but not the menu registry}
\date{September 2026}
\maketitle

This is a fidelity control for \texttt{menu\_lifecycle.tex}. Only
\texttt{DepartLapsed} differs; every other schema is copied verbatim.

\begin{zed}
  [CONN]
\end{zed}

\begin{zed}
  DSTATE ::= dnone | ddelivered | dgone
\end{zed}

\begin{zed}
  FSTATE ::= fclear | fset
\end{zed}

\begin{zed}
  APHASE ::= aidle | astore
\end{zed}

\begin{zed}
  OPHASE ::= oidle | oput
\end{zed}

\begin{schema}{MenuReg}
  registered : \power CONN \\
  identified : \power CONN \\
  lapsed : \power CONN \\
  hasInbox : \power CONN \\
  menuOwner : \power CONN \\
  reported : DSTATE \\
  lost : FSTATE \\
  admitPhase : APHASE \\
  admitFor : \power CONN \\
  offerPhase : OPHASE \\
  offerFor : \power CONN
\where
  identified \subseteq registered \\
  lapsed \subseteq registered \\
  \# admitFor \leq 1 \\
  \# offerFor \leq 1
\end{schema}

\begin{schema}{Init}
  MenuReg'
\where
  registered' = \emptyset \\
  identified' = \emptyset \\
  lapsed' = \emptyset \\
  hasInbox' = \emptyset \\
  menuOwner' = \emptyset \\
  reported' = dnone \\
  lost' = fclear \\
  admitPhase' = aidle \\
  admitFor' = \emptyset \\
  offerPhase' = oidle \\
  offerFor' = \emptyset
\end{schema}

\begin{schema}{Connect}
  \Delta MenuReg \\
  c? : CONN
\where
  registered' = registered \cup \{ c? \} \\
  lapsed' = lapsed \setminus \{ c? \} \\
  identified' = identified \\
  hasInbox' = hasInbox \\
  menuOwner' = menuOwner \\
  reported' = reported \\
  lost' = lost \\
  admitPhase' = admitPhase \\
  admitFor' = admitFor \\
  offerPhase' = offerPhase \\
  offerFor' = offerFor
\end{schema}

\begin{schema}{Identify}
  \Delta MenuReg \\
  c? : CONN
\where
  c? \in registered \\
  identified' = identified \cup \{ c? \} \\
  registered' = registered \\
  lapsed' = lapsed \\
  hasInbox' = hasInbox \\
  menuOwner' = menuOwner \\
  reported' = reported \\
  lost' = lost \\
  admitPhase' = admitPhase \\
  admitFor' = admitFor \\
  offerPhase' = offerPhase \\
  offerFor' = offerFor
\end{schema}

\begin{schema}{Arm}
  \Delta MenuReg \\
  c? : CONN
\where
  c? \in registered \\
  hasInbox' = hasInbox \cup \{ c? \} \\
  registered' = registered \\
  identified' = identified \\
  lapsed' = lapsed \\
  menuOwner' = menuOwner \\
  reported' = reported \\
  lost' = lost \\
  admitPhase' = admitPhase \\
  admitFor' = admitFor \\
  offerPhase' = offerPhase \\
  offerFor' = offerFor
\end{schema}

\begin{schema}{LeaseLapse}
  \Delta MenuReg \\
  c? : CONN
\where
  c? \in registered \\
  c? \notin lapsed \\
  lapsed' = lapsed \cup \{ c? \} \\
  registered' = registered \\
  identified' = identified \\
  hasInbox' = hasInbox \\
  menuOwner' = menuOwner \\
  reported' = reported \\
  lost' = lost \\
  admitPhase' = admitPhase \\
  admitFor' = admitFor \\
  offerPhase' = offerPhase \\
  offerFor' = offerFor
\end{schema}

\begin{schema}{SetMenu}
  \Delta MenuReg \\
  c? : CONN
\where
  c? \in registered \\
  c? \in identified \\
  lapsed' = lapsed \setminus \{ c? \} \\
  hasInbox' = hasInbox \cup \{ c? \} \\
  menuOwner' = menuOwner \cup \{ c? \} \\
  registered' = registered \\
  identified' = identified \\
  reported' = reported \\
  lost' = lost \\
  admitPhase' = admitPhase \\
  admitFor' = admitFor \\
  offerPhase' = offerPhase \\
  offerFor' = offerFor
\end{schema}

\begin{schema}{SetMenuRefused}
  \Delta MenuReg \\
  c? : CONN
\where
  c? \notin identified \\
  lapsed' = lapsed \setminus \{ c? \} \\
  registered' = registered \\
  identified' = identified \\
  hasInbox' = hasInbox \\
  menuOwner' = menuOwner \\
  reported' = reported \\
  lost' = lost \\
  admitPhase' = admitPhase \\
  admitFor' = admitFor \\
  offerPhase' = offerPhase \\
  offerFor' = offerFor
\end{schema}

\begin{schema}{DepartGraceful}
  \Delta MenuReg \\
  c? : CONN
\where
  c? \in registered \\
  registered' = registered \setminus \{ c? \} \\
  identified' = identified \setminus \{ c? \} \\
  lapsed' = lapsed \setminus \{ c? \} \\
  hasInbox' = hasInbox \setminus \{ c? \} \\
  menuOwner' = menuOwner \setminus \{ c? \} \\
  reported' = reported \\
  lost' = lost \\
  admitPhase' = admitPhase \\
  admitFor' = admitFor \\
  offerPhase' = offerPhase \\
  offerFor' = offerFor
\end{schema}

% -- THE DEFECT: the lapse path reaps the inbox but not the menu registry.
% menuOwner is carried UNCHANGED, so a timer-reaped owner keeps its bar.
\begin{schema}{DepartLapsed}
  \Delta MenuReg \\
  c? : CONN
\where
  c? \in registered \\
  c? \in lapsed \\
  registered' = registered \setminus \{ c? \} \\
  identified' = identified \setminus \{ c? \} \\
  lapsed' = lapsed \setminus \{ c? \} \\
  hasInbox' = hasInbox \setminus \{ c? \} \\
  menuOwner' = menuOwner \\
  reported' = reported \\
  lost' = lost \\
  admitPhase' = admitPhase \\
  admitFor' = admitFor \\
  offerPhase' = offerPhase \\
  offerFor' = offerFor
\end{schema}

\begin{schema}{MenuDeliver}
  \Delta MenuReg \\
  c? : CONN
\where
  c? \in registered \\
  c? \in hasInbox \\
  reported' = ddelivered \\
  registered' = registered \\
  identified' = identified \\
  lapsed' = lapsed \\
  hasInbox' = hasInbox \\
  menuOwner' = menuOwner \\
  lost' = lost \\
  admitPhase' = admitPhase \\
  admitFor' = admitFor \\
  offerPhase' = offerPhase \\
  offerFor' = offerFor
\end{schema}

\begin{schema}{MenuProviderGone}
  \Delta MenuReg \\
  c? : CONN
\where
  \lnot (c? \in registered \land c? \in hasInbox) \\
  reported' = dgone \\
  registered' = registered \\
  identified' = identified \\
  lapsed' = lapsed \\
  hasInbox' = hasInbox \\
  menuOwner' = menuOwner \\
  lost' = lost \\
  admitPhase' = admitPhase \\
  admitFor' = admitFor \\
  offerPhase' = offerPhase \\
  offerFor' = offerFor
\end{schema}

\end{document}
